%% kapa: exam=1 round=27 year=2016 period=1 form=A subject=civil_law subject_ko=민법 seq=24 number=24
%% source: Q-Net 감정평가사 기출문제 (cst00309) artlSeq=5169402 file=extracted/1008/5169402/제27회 감정평가사 1차 문제/제27회 감정평가사 1교시문제(A형).pdf
%% conversion: kordoc 4.24.51 -> tools/gen1.py (auto); review=proofread-sonnet; flags=none
\begin{question}{24}
등기의 효력에 관한 설명으로 옳은 것은? (다툼이 있으면 판례에 따름)
\begin{choices}
  \item 소유권이전청구권 보전을 위한 가등기가 있어도 소유권이전등기를 청구할 어떤 법률관계가 있다고 추정되지 않는다.
  \item 허무인(虛無人)으로부터 이어받은 소유권이전등기의 경우에도 그 등기명의자의 소유권은 추정된다.
  \item 신축된 건물의 소유권보존등기 명의자는 실제로 그 건물을 신축한 자가 아니더라도 적법한 권리자로 추정된다.
  \item 등기가 원인 없이 말소된 경우, 그 회복등기가 마쳐지기 전이라면 말소된 등기의 등기명의인은 적법한 권리자로 추정되지 않는다.
  \item 소유권이전등기 명의자는 그 전(前) 소유자에 대하여 적법한 등기원인에 의해 소유권을 취득한 것으로 추정되지 않는다.
\end{choices}
\end{question}
\answer{1}
